% CVPR 2026 Paper Template; see https://github.com/cvpr-org/author-kit

\documentclass[10pt,twocolumn,letterpaper]{article}

%%%%%%%%% PAPER TYPE  - PLEASE UPDATE FOR FINAL VERSION
% \usepackage{cvpr}              % To produce the CAMERA-READY version
\usepackage[review]{cvpr}      % To produce the REVIEW version
% \usepackage[pagenumbers]{cvpr} % To force page numbers, e.g. for an arXiv version

% Import additional packages in the preamble file, before hyperref
\input{preamble}

% It is strongly recommended to use hyperref, especially for the review version.
% hyperref with option pagebackref eases the reviewers' job.
% Please disable hyperref *only* if you encounter grave issues, 
% e.g. with the file validation for the camera-ready version.
%
% If you comment hyperref and then uncomment it, you should delete *.aux before re-running LaTeX.
% (Or just hit 'q' on the first LaTeX run, let it finish, and you should be clear).
\definecolor{cvprblue}{rgb}{0.21,0.49,0.74}
\usepackage[pagebackref,breaklinks,colorlinks,allcolors=cvprblue]{hyperref}
\usepackage{multirow}

%%%%%%%%% PAPER ID  - PLEASE UPDATE
\def\paperID{16} % *** Enter the Paper ID here
\def\confName{CVPR}
\def\confYear{2026}

%%%%%%%%% TITLE - PLEASE UPDATE
\title{FiZK: Federated Integrity with Zero Knowledge\\
Hierarchical FL with ZK-Firewall Defense and IVC Proof Folding}

%%%%%%%%% AUTHORS - PLEASE UPDATE
\author{a1\\
Computer Engineering\\
St.\ Vincent Pallotti College of Engineering and Technology\\
Nagpur, India\\
{\tt\small a1@gmail.com}
\and
a3\\
Computer Engineering\\
St.\ Vincent Pallotti College of Engineering and Technology\\
Nagpur, India\\
{\tt\small a3@gmail.com}
\and
a2\\
Computer Engineering\\
St.\ Vincent Pallotti College of Engineering and Technology\\
Nagpur, India\\
{\tt\small a2@gmail.com}
\and
a4\\
Computer Engineering\\
St.\ Vincent Pallotti College of Engineering and Technology\\
Nagpur, India\\
{\tt\small a4@gmail.com}
\and
a5\\
Computer Engineering\\
St.\ Vincent Pallotti College of Engineering and Technology\\
Nagpur, India\\
{\tt\small a5@gmail.com}
}

\begin{document}
\maketitle

% ============================================================
% ABSTRACT
% ============================================================
\begin{abstract}

\itshape
Federated learning enables a collaborative learning approach without requiring data sharing; however, it also makes it vulnerable to various attacks, especially Byzantine attacks, due to its distributed learning approach. Byzantine attacks are a major threat to federated learning systems, as they can compromise the integrity of the models by incorporating malicious gradients that perform data poisoning attacks. Although robust aggregation makes federated learning robust against Byzantine attacks, it is not guaranteed that the updates from clients are correct. In this paper, we propose a novel hierarchical federated learning approach called FiZK that utilizes various features of zero-knowledge gradient validation, commitment, and robust aggregation to ensure verifiable integrity within the federated learning paradigm. The proposed FiZK approach includes a ZK gradient validation layer that validates gradient norms during aggregation. Additionally, it utilizes a commitment approach that includes a Merkle tree and proof folding to ensure efficiency during Byzantine attack handling. The efficiency of FiZK is demonstrated by evaluating its performance during various attack scenarios while ensuring robustness to guarantee accurate performance during Byzantine attacks.
\end{abstract}


% ============================================================
\section{Introduction}
\label{sec:intro}
% ============================================================

\subsection{Motivation}

Federated Learning (FL) is a framework used for collaborative learning without sharing sensitive raw data \cite{li2020federated,kairouz2021advances}. This is the major reason FL is being adopted in privacy-oriented industries such as the health sector, finance sector, and mobile computing. However, the training process is vulnerable to security attacks whereby a malicious client may send modified gradients that may reduce the accuracy of the global model.

The types of attacks vary. The attacks may also be connected. Byzantine clients may use poisoned gradients. This may lead to a decrease in the accuracy of the global model. The poisoned gradients may also lead to the insertion of backdoors in the global model \cite{blanchard2017machine,fang2020local}. The gradients may be dropped by the compromised aggregator. This type of attack is called the \emph{aggregator trust problem}. In FL, the gradients sent by the clients are not cryptographically linked. Hence, it is essentially impossible to trace the source of the attack. The computation overhead of the Multi-Krum method is also quite high, $O(n^2)$ which is a major drawback of the method. Most of the existing literature deals with just one type of attack. Most of the robust aggregation techniques proposed in the literature, including the Krum method \cite{blanchard2017machine} and the Trimmed Mean method \cite{yin2018byzantine}, consider the case where only a small percentage of clients may be sending poisoned gradients. However, the client sending the gradients may also be dishonestly. On the contrary, the secure aggregation method proposed in the literature \cite{bonawitz2017practical} uses the secret sharing method. The method ensures the privacy of the gradients. However, the method is also vulnerable to Byzantine attacks. The FLTrust method proposed in the literature \cite{cao2021fltrust} assumes that the server has a clean root dataset. However, the method violates the data minimization principle.

\subsection{Key Insight}
% : ZK-First, Fold-Then-Verify}

% We present FiZK, which addresses these shortcomings with two complementary cryptographic mechanisms. On every client, a zero-knowledge proof attests that the gradient sum is correct and that $\|\nabla w_i\|_2 \leq B$ (with $B$ derived from the server-side norm distribution using $\mathrm{median} \pm k \cdot \mathrm{MAD}$). This ZK-Firewall places a hard safety envelope around each gradient \emph{before} it reaches aggregation, so high-magnitude attacks are blocked at the proof-validation stage instead of being detected only after damage is done. At the galaxy level, the aggregator folds each client's proof into one compact proof $\pi_g^{\mathsf{fold}}$ using IVC-based folding through the ProtoGalaxy scheme \cite{eagen2023protogalaxy}, as implemented in the Sonobe library \cite{sonobe2024} on BN254/Grumpkin curves.

\textbf{FiZK} combines two cryptographic mechanisms. First, each client submits a zero-knowledge proof attesting that its per-layer gradient sums are correctly computed and within a server-derived norm bound $B_l$ (anchored at the 25th percentile of client norms, keeping it in the honest regime until $\alpha \approx 0.75$). A gradient that violates the bound cannot satisfy the ZK circuit, so the client is rejected \emph{cryptographically non-repudiably} this is the \emph{ZK-Firewall} (formally defined in \cref{sec:zkp}). Second, clients are grouped into clusters called \emph{galaxies}; the galaxy aggregator folds all per-client proofs into a single compact proof $\pi_g^{\mathsf{fold}}$ using ProtoGalaxy IVC~\cite{eagen2023protogalaxy} via the Sonobe library~\cite{sonobe2024}, reducing server-side verification to $O(1)$ per galaxy.

\subsection{Our contribution}

% We efforts are focused towards:

% \begin{enumerate}
% \item \textbf{Four-phase verifiable FL pipeline}: A formally specified
%   commit-validate-fold protocol (Phase~1: client commitment and ZKP
%   generation; Phase~2: Merkle and ZK verification; Phase~3: Multi-Krum
%   selection and galaxy-level IVC folding; Phase~4: global Merkle
%   verification and model update).

% \item \textbf{ZK-Firewall}: A zero-knowledge gate proving gradient sum
%   correctness and norm boundedness ($\|\nabla w_i\|_2 \leq B$) via a
%   2-constraint arithmetic circuit, enforced before any gradient enters
%   aggregation.

% \item \textbf{Galaxy-Level IVC Folding}: Application of the ProtoGalaxy
%   IVC folding scheme (via Sonobe on BN254/Grumpkin) to compress
%   per-client proofs within each galaxy into a single folded proof
%   $\pi_g^{\mathsf{fold}}$.

% \item \textbf{Merkle Integrity Lock}: A two-level hierarchical Merkle
%   tree providing $O(\log n)$ gradient commitment verification and an
%   immutable forensic audit trail linking every client submission to a
%   published root.

% \item \textbf{Empirical evaluation}: Comparison across adaptive,
%   label-flip, and Gaussian noise attacks on MNIST at 10-client scale,
%   demonstrating that ZK-based filtering extends effective Byzantine
%   tolerance beyond Multi-Krum's theoretical bound of $f < n/2 - 2$.
% \end{enumerate}

% In our work, we propose \textbf{FiZK}, a hierarchical federated learning framework that targets the integrity and robustness of Byzantine robust distributed training. Our first contribution is the \emph{ZK-Firewall}, in which a zero-knowledge proof is created by each client to verify the correctness of the sum of gradients and the norm of gradients. Our second contribution is a scalable verification method that groups clients into \emph{galaxies} and leverages IVC folding via ProtoGalaxy to condense all client proofs into a single proof. Our third contribution is the integration of a two-level Merkle commitment scheme to store all client submissions. Our fourth contribution is a thorough experimental evaluation of our algorithm on linear classifiers (MNIST) and convolutional neural networks (\textbf{LeNet-5 trained from scratch on Fashion-MNIST}) to show its efficacy against various types of Byzantine attacks such as model poisoning, label flip, backdoor, and adaptive attacks. On MNIST linear models, our algorithm maintains accuracy over 92\% under model poisoning with $\alpha=0.60$, whereas Multi-Krum collapses to near-random accuracy ($<12\%$) under $\alpha\geq0.50$. On Fashion-MNIST CNNs, FiZK achieves perfect detection (TPR\,=\,100\%, FPR\,=\,0\%, F1\,=\,1.0) across all attack types at $\alpha=0.30$, maintaining 81.4--81.5\% accuracy. Multi-Krum is highly unstable under poisoning (32.5\,$\pm$\,29.3\% under IID model poisoning at $\alpha=0.30$) and degrades completely at $\alpha\geq0.50$.
In our work, we propose a hierarchical federated learning framework called \textbf{FiZK} with a focus on the integrity and robustness of Byzantine robust distributed training. Our first contribution is the development of the \textit{ZK-Firewall} method in which a zero-knowledge proof is generated by each client to prove the correctness of the sum of gradients and the norm of gradients. Our second contribution is a scalable verification method in which clients are grouped in a hierarchical manner called \textit{galaxies} and use IVC folding through the use of the ProtoGalaxy method. Our third contribution is the incorporation of a two-level Merkle commitment scheme in which all client submissions are stored. Our fourth contribution is an extensive experimental evaluation of the algorithm using linear models (MNIST) and CNNs (\textbf{LeNet-5} trained from scratch using FashionMNIST) in defending against various types of Byzantine attacks: model poisoning, label flip attack, backdoor attack, and adaptive attack. Our algorithm maintains an accuracy of over 92\% in defending against model poisoning attacks with $\alpha=0.60$, whereas Multi-Krum collapses with accuracy of less than 12\% with $\alpha\geq0.50$. Our algorithm also achieves a perfect score in defending against all types of Byzantine attacks with TPR=100\%, FPR=0\%, and F1-score of 1.0 on CNNs with $\alpha=0.30$, whereas Multi-Krum is unstable with accuracy of only 32.5$\pm$29.3\% with IID model poisoning with $\alpha=0.30$. Our algorithm maintains an accuracy of 81.4-81.5\% with $\alpha=0.30$, whereas Multi-Krum collapses with $\alpha\geq0.50$.

\section{Related work}
\label{sec:related}
% Most work on robust Federated Learning has focused either on statistical defenses or on cryptographic privacy; verifiable integrity at the level of individual client contributions is comparatively unexplored. Byzantine-robust aggregation rules like Krum \cite{blanchard2017machine}, coordinate-wise median \cite{yin2018byzantine}, and FLTrust \cite{cao2021fltrust} provide useful resilience or root-dataset-based heuristic, but they cannot supply pre-revelation gradient norm bounding or produce compact audit proofs. Blockchain-based FL \cite{kim2020blockchained} does offer global verifiability, yet the associated latency overhead is often impractical. Existing zero-knowledge proposals for FL have largely targeted server-side aggregation correctness rather than per-client gradient properties, and the broader zkML community has so far concentrated on verifiable inference.

% FiZK fills the gap between these lines of work. We use the ProtoGalaxy IVC scheme \cite{eagen2023protogalaxy}, accessed through the Sonobe library \cite{sonobe2024}, to build a lightweight two-constraint arithmetic circuit that checks gradient sum correctness and norm bounding. The result complements secure aggregation protocols like \cite{bonawitz2017practical}, which protect privacy through secret sharing but do not, on their own, provide the integrity guarantees or proof-compression efficiency that our architecture delivers.

% Most of the solid Federated Learning research uses statistical or cryptographic approaches, but the verifiability of client integrity has not been extensively explored. Approaches such as Krum~\cite{blanchard2017machine}, coordinate-wise median~\cite{yin2018byzantine}, and FLTrust~\cite{cao2021fltrust}  are attempts at mitigating malicious clients or using a root dataset, but they do not include the gradient norm constraints before aggregation.

% Another direction in this area is system-level transparency using a blockchain-based federated learning approach \cite{kim2020blockchained}, where global verifiability is provided, although this comes with a high latency and scalability cost. Concurrently, zero-knowledge techniques are starting to emerge in ML systems, although most recent work in zkML focuses on verifiable inference, not training. Zero-knowledge techniques proposed in existing work on federated learning verify the correctness of a server-side aggregation mechanism, not client gradients.
 Most of the existing Federated Learning literature on solid research is based on either statistical or cryptographic methods, but the integrity of clients has not been extensively verified. Krum~\cite{blanchard2017machine}, coordinate-wise median~\cite{yin2018byzantine}, FLTrust~\cite{cao2021fltrust} are some attempts to address malicious clients or use a root dataset, but they do not include gradient norm constraints before aggregation.
 
Another direction in this space is system-level transparency through a blockchain-based federated learning approach \cite{kim2020blockchained}, where global verifiability is achieved, albeit with a high cost in terms of latency and scalability. At the same time, zero knowledge techniques are beginning to emerge in ML systems, although much of the recent work in zkML has been on verifiable inference, not training. Zero knowledge techniques that are proposed in existing work on federated learning verify correctness for a server-side aggregation mechanism, not client gradients.

% ============================================================
\section{Proposed approach}
\label{sec:approach}
% ============================================================

\subsection{Threat model}
\label{sec:threat}

% We consider a strong adaptive adversary that attacks concurrently at three levels. First, at the client level, up to $\alpha = 60\%$ of the clients can behave as Byzantine adversaries and may submit any gradients or strategically manipulate gradients to the aggregator, including label flipping, backdoors injection, model poisoning, or structured noise. Second, at the aggregator level, the aggregator can manipulate or tamper with the gradients or client identities. Finally, at the network level, an adversary can intercept or manipulate gradient messages.
We consider a strong adaptive adversary that attacks simultaneously at three levels. First, at the client level, up to $\alpha = 60\%$ of the clients behave as Byzantine attackers and may report any gradients or strategically manipulate gradients to the aggregator, including label flipping, backdoor injection, model poisoning, and structured noise. Second, at the aggregator level, the aggregator may manipulate or tamper with the gradients or client identities. Third, at the network level, an adversary may intercept or manipulate the gradient messages.

We make the following assumptions:
(i) the local hardware of the clients is trusted;
(ii) SHA-256 and the ZK argument system used are computationally sound;
(iii) Messages can be delivered within a bounded time;
(iv) Every client has a unique registered key pair.

% ============================================================
\subsection{System overview}
\label{sec:arch}
% ============================================================

% FiZK has a two-level hierarchy for the participants, as depicted in \cref{fig:arch}. In the first level, the clients are divided into groups called \textbf{galaxies}. Each galaxy has a corresponding \textbf{Galaxy Aggregator} responsible for the collection of local commitments, zero knowledge verification, Byzantine aggregation, and proof compression. In the second level, the \textbf{Global Server} is responsible for collecting the verified results from each galaxy, conducting the cross-galaxy anomaly test, and the application of the update to the model.

% In more technical terms, let $n$ denote the total number of clients. Let $G$ denote the number of galaxies. Then, the clients are divided into $G$ disjoint sets $\mathcal{C}_1, \mathcal{C}_2, \dots, \mathcal{C}_G$. Thus, $\bigcup_{g=1}^{G} \mathcal{C}_g = \mathcal{C}$. For each galaxy $g$, the corresponding \textbf{Galaxy Aggregator} computes the aggregated weights $\bar{w}_g$. It also computes the Merkle root $R_g$ for all client submissions. Furthermore, the \textbf{Galaxy Aggregator} computes the compact proof $\pi_g^{\mathrm{fold}}$. This proof is for the validity of all client computations without the need for the global server to reverify the proofs.
The hierarchy of the participants is two-levelled in the case of FiZK. This is illustrated in \cref{fig:arch}. For the first level, the clients are partitioned into groups called \textbf{galaxies}. For each galaxy, a \textbf{Galaxy Aggregator} is responsible for the collection of local commitments, zero knowledge verification, Byzantine aggregation, and proof compression. For the second level, the \textbf{Global Server} is responsible for the collection of verified results from all the galaxies, cross-galaxy anomaly test, and the application of the update to the model.

More specifically, let $n$ be the total number of clients. Let $G$ be the total number of galaxies. The clients are partitioned into disjoint sets $\mathcal{C}_1$, $\mathcal{C}_2$, ..., $\mathcal{C}_G$. Hence, $\bigcup_{g=1}^{G} \mathcal{C}_g = \mathcal{C}$. For each galaxy $g$, the Galaxy Aggregator computes the aggregated weights $\bar{w}_g$. It also computes the Merkle root $R_g$ of all client submissions. Further, the Galaxy Aggregator computes the proof $\pi_g^{\mathrm{fold}}$. It is a proof of validity of all client computations without the need for the global server to reverify the proofs.

\begin{figure*}[t]
    \centering
    \includegraphics[width=0.8\textwidth]{new architecture.png}
    \caption{The FiZK architecture. \textit{Edge Clients} generate cryptographic artifacts in Phase~1 and reveal gradients through the \textit{Layer~1 Crypto Integrity} gate in Phase~2. The \textit{Filters \& Defense} module applies statistical screening and reputation-based filtering (Phases~2--3), while the \textit{Aggregation} module compresses per-client proofs into a single galaxy proof $\pi_g^{\mathsf{fold}}$ via ProtoGalaxy IVC folding. The \textit{Global Server} performs Merkle verification, cross-galaxy anomaly detection (Layer~5), and the final model update (Phase~4).}
    \label{fig:arch}
\end{figure*}

% ============================================================
\subsection{Four-phase training protocol}
\label{sec:protocol}
% ============================================================

Each training round goes through four phases shown in \cref{fig:arch}. The first two phases implement a \emph{commit-then-reveal} protocol: clients lock themselves to a gradient before exposing it, so any attempt at late-stage manipulation is cryptographically detectable.

\paragraph{Phase 1  Commitment and ZKP Generation}
Each client trains its model and calculates the gradients $\nabla w_i^{(r)}$. 
before sending any information to the Galaxy Aggregator, each client generates two items (the Edge Clients block in \cref{fig:arch}).

These items include the binding commitment
\begin{equation}
h_i = H(\nabla w_i^{(r)} \,\|\, r \,\|\, \text{nonce}),
\label{eq:commitment}
\end{equation}
where the gradient is hidden but locked to a specific value.

The zero-knowledge proof $\pi_i^{(\mathrm{zk})}$ is used to show that the gradients have been correctly summed. Moreover, the norm of the gradients is verified.

The per-layer norm bound is computed as
\begin{equation}
B_l = Q_{0.25}\!\left(\left\{\|\nabla w_i\|_2^{(l)}\right\}_{i=1}^{n}\right) \times (1 + \gamma),\quad \gamma = 5.
\label{eq:norm_bound}
\end{equation}

Anchoring at the 25th percentile keeps $B_l$ in the honest regime for $\alpha < 0.75$.

The client sends $h_i$ and $\pi_i^{(\mathrm{zk})}$ to the Galaxy Aggregator. 
The aggregator forms a galaxy Merkle tree from all the hash values. Finally, the immutable root $R_g^{(r)}$ is published.

\paragraph{Phase 2: Gradient Revelation and ZK-Firewall}
After the locking of the gradients in the first phase, the gradients are revealed. A two-stage cryptographic gate is imposed by the \textit{Module: Layer~1 Crypto Integrity} (\cref{fig:arch}).

The \textbf{Merkle Lock} checks whether the hash of the revealed gradients equals the published root $R_g^{(r)}$. If the hash of the revealed gradients does not equal the published root because of tampering or the presence of a man-in-the-middle attacker, the gradients are rejected.

In the \textbf{ZK Firewall}, the correctness of the gradients is verified. If the norm of the gradient falls outside the permissible range B, it is rejected.
% Together, these two checks constitute a hard cryptographic safety envelope around every submitted update.


% \textbf{Claim 3.1 (Binding Commitment):} Suppose SHA-256 is still collision-free. Then, once the value of $R_g$ is revealed, it is impossible for an adversary to provide a gradient $w'$  with a different gradient from $\nabla w' \neq \nabla w_i$, such that the hash $h_i$ is the same.
\paragraph{Phase 3  Multi-Layer Defense and Proof Folding}
Gradient data is received and goes through the Filters \& Defense and Aggregation blocks as shown in \cref{fig:arch}.
On the defense side, a statistical test is implemented in Layer 2. The test is carried out on gradients from four different tests: norm, cosine direction, coordinate distribution, and KL divergence from the group. Gradients are rejected if they fail less than two out of four tests.
In Layer 3, the gradients are filtered using the multi-krum method to select gradients that are closest to each other and then discard those gradients that are farther away from the group before computing the galaxy aggregate, represented as $\bar{w}_g$. In Layer 4, the client reputation score $\mathcal{R}_i$ is exponentially updated as follows:
\begin{equation}
\label{eq:reputation}
\mathcal{R}_i(t{+}1) = (1-\lambda)\,\mathcal{R}_i(t) + \lambda\, B_i(t), \quad \lambda = 0.1,
\end{equation}
where $\lambda$ is the learning rate and is fixed at 0.1, and $B_i(t)$ is the cumulative sum of pass/fail indicators.

On the proof side, the aggregation block uses ProtoGalaxy IVC Folding. It folds the proof as follows:
~\cite{eagen2023protogalaxy}

\begin{equation}
\label{eq:ivc_fold_def}
\pi_g^{\mathsf{fold}} = \mathcal{F}_{\mathsf{IVC}}\!\left(\dots\mathcal{F}_{\mathsf{IVC}}(\pi_1^{zk}, \pi_2^{zk})\dots, \pi_m^{zk}\right)
\end{equation}

As the size of the proof is independent of the value of $m$, the verification time at the Global Server is constant time complexity $O(1)$. Then the Galaxy Aggregator sends the triplet $(\bar{w}_g,\, \pi_g^{\mathsf{fold}},\, R_g)$  to the Global Server.

\paragraph{Phase 4  Global Verification and Model Update}

This Global Server (\cref{fig:arch}) processes all the galaxies through a three-stage pipeline. First, the Merkle Verifier gathers all the galaxy roots to construct the global tree and verifies each of the $R_g$ values against the global root $R$. Subsequently, the Layer 5 Galaxy Defense checks the aggregate updates of the galaxies as a whole to cross-check for anomalous galaxies by checking norm deviation and similarities between the galaxies. Finally, the Global Defense Pipeline aggregates the good updates to calculate the new weights $w^{(r+1)}$ for the entire participant set.


% ============================================================
\subsection{ZK-Firewall and IVC folding}
\label{sec:zkp}
% ============================================================

The architecture's trustworthiness rests on two cryptographic mechanisms: the per-client ZK circuit and the galaxy-level IVC folding scheme.

\subsubsection{Client zero-knowledge proof}

Each client proves, without revealing its gradient, that it belongs to the relation:
\begin{equation}
\mathcal{R}_{\mathrm{zk}} = \left\{
(\nabla w_i,\, x_{\mathrm{pub}}) :\!
\begin{array}{l}
  \textstyle\sum\nolimits_j (\nabla w_i)_j^{(l)} = S_i^{(l)} \;\forall l \\[3pt]
  \bigl(S_i^{(l)}\bigr)^2 \leq B_l^2 \;\forall l
\end{array}
\right\}
\label{eq:zk_relation}
\end{equation}
where $S_i^{(l)}$ is the scalar sum of all parameters in layer $l$ of client $i$'s gradient, and $B_l$ is the server-side per-layer bound computed before the round as:
\begin{equation}
B_l = Q_{0.25}\!\left(\{\|\nabla w_i\|_2^{(l)}\}_{i=1}^n\right) \times (1 + \gamma),\quad \gamma = 5
\label{eq:bound}
\end{equation}
where $Q_{0.25}$ refers to the 25th percentile of client norms for the layer $l$. The two conditions in $\mathcal{R}_\mathrm{zk}$ are referred to as \emph{sum correctness} (the sum equals what was committed) and \emph{per-layer bound} (the sum does not exceed $B_l$).

To verify $\bigl(S_i^{(l)}\bigr)^2 \leq B_l^2$, the prover must express the difference $B_l^2 - \bigl(S_i^{(l)}\bigr)^2$ as a non-negative 64-bit integer. If the gradient sum is too large, the difference will be negative and cannot be expressed in this form; proof generation fails and the client will be rejected without a valid certificate. Proofs are generated using BN254/Grumpkin curves with the Sonobe library~\cite{sonobe2024}.


\subsubsection{Galaxy-level IVC folding}

After the defense pipeline has selected $m$ clients, the aggregator folds these clients' proofs in sequence, using Incrementally Verifiable Computation (IVC)~\cite{valiant2008incrementally}, beginning from $\mathrm{acc}_1 = \pi_1^{zk}$:
\begin{equation}
\mathrm{acc}_{k} = \mathcal{F}_{\mathsf{IVC}}(\mathrm{acc}_{k-1},\, \pi_k^{zk}), \quad k = 2, \dots, m.
\label{eq:ivc_fold}
\end{equation}
The final accumulator $\pi_g^{\mathsf{fold}} = \mathrm{acc}_m$ is constant size independent of $m$. Galaxy IVC folding reduces $m$ per-client proof verifications to $O(1)$. Combined with the two-level Merkle tree, total server-side verification cost per round is $O(G \log n)$, where $G$ is the number of galaxies.

% ============================================================
\section{Experiments and results}
\label{sec:eval}
% ============================================================

%                     
\subsection{Implementation and setup}
%                     

\subsubsection{Implementation}

The implementation uses \textbf{PyTorch}~\cite{paszke2019pytorch} for model training, SHA-256~\cite{nist2015sha} for Merkle trees, and the \textbf{Sonobe IVC library}~\cite{sonobe2024} (ProtoGalaxy on \textbf{BN254/Grumpkin} curves) via a Rust-Python bridge (\textbf{Maturin}~\cite{maturin2024}). Defense mechanisms include Multi-Krum~\cite{blanchard2017machine}, coordinate-wise median, trimmed mean~\cite{yin2018byzantine}, statistical anomaly detection (norm analysis, cosine similarity, MAD, KL-divergence), and EWMA reputation scoring.

\subsubsection{Common experimental configuration}

\textbf{Baselines}: Vanilla FedAvg~\cite{mcmahan2017communication} (no defense), Multi-Krum~\cite{blanchard2017machine}, and FiZK (ZKP~+~Merkle~+~IVC folding).

\textbf{Scale}: 10--50 clients in 2--10 galaxies via round-robin assignment.

\textbf{Byzantine fractions}: $\alpha \in \{0.30, 0.50, 0.60\}$.

\textbf{ZK bound}: $B_l = Q_{0.25}(\{\|\nabla w_i\|_2^{(l)}\}) \times 6$ per layer.

\textbf{Attacks}: Model poisoning ($\|\nabla w_i\|_2 \gg B$)~\cite{fang2020local}, label flipping~\cite{tolpegin2020data}, backdoor (pixel-pattern trigger)~\cite{gu2019badnets}, adaptive (gradient-aware evasion), gradient substitution, and compromised aggregator.

\textbf{Trials}: 3 independent trials per configuration (seed varied).

%                     
\subsection{Experiment 1: MNIST with linear classifier}
\label{sec:exp_mnist}
%                     

\subsubsection{Setup}

\textbf{Dataset}: MNIST~\cite{lecun1998gradient} (60K train, 10K test); IID and Dirichlet non-IID ($\beta=0.5$) partitioning.

\textbf{Model}: Single fully-connected layer(7850), trained with cross-entropy loss.

\textbf{Training}: 10 rounds; 1 local epoch; batch size 64; SGD lr=0.01.

\subsubsection{Results}

\cref{fig:mnist_comparison} presents accuracy for all defenses across IID and non-IID scenarios at $\alpha\in\{0.30,0.50\}$. Under IID model poisoning, FiZK holds at 92.0-92.1\% at both Byzantine fractions. Multi-Krum collapses to $\leq$11\% at $\alpha=0.50$ (near-random) and Vanilla collapses ($\leq$6.3\%). Under Dirichlet non-IID scenarios, FiZK holds at 87.3-89.7\% at $\alpha=0.30$ for all attacks. At $\alpha=0.50$, FiZK holds at 88.7-89.7\% while Multi-Krum degrades to 8.4\% under label flip and model poisoning. The effect of backdoor attacks is minimal for all defenses ($\geq$83\%). Under adaptive attack (\cref{tab:accuracy}), FiZK holds at 92.06\% while Multi-Krum holds at 90.85\% at $\alpha=0.50$. The computational cost is $8.80\times$ (\cref{tab:overhead}), mainly due to CPU-intensive ZKP computations that account for 75.8\% of round time.

\textbf{Key finding}: FiZK extends Byzantine tolerance well beyond Multi-Krum's theoretical limit ($f<n/2-2$, where $f=\lfloor\alpha n\rfloor$ is the Byzantine client count), maintaining $\geq$92\% IID accuracy even under model poisoning at $\alpha=0.60$, by anchoring $B_l$ at $Q_{0.25}$ which remains in the honest regime until $\alpha\approx0.75$.

\begin{figure*}[t]
  \centering
  \includegraphics[width=\textwidth]{mnist_full_comparison.pdf}
  \caption{MNIST attack robustness: test accuracy (\%) for Vanilla FedAvg, Multi-Krum, and FiZK under Backdoor, Label Flip, and Model Poisoning. Top row: IID; bottom row: Dirichlet non-IID ($\beta=0.5$). Byzantine fractions $\alpha\in\{0.30,0.50\}$. FiZK consistently holds $\geq$87\% under non-IID and $\geq$92\% IID under all attacks; Multi-Krum collapses at $\alpha=0.50$ under label flip and model poisoning.}
  \label{fig:mnist_comparison}
\end{figure*}

%                     
\subsection{Experiment 2: Fashion-MNIST with LeNet-5 CNN}
\label{sec:exp_cnn}
%                     

\subsubsection{Setup}

\textbf{Dataset}: Fashion-MNIST~\cite{xiao2017fashion} (60K train, 10K test); IID and Dirichlet non-IID.

\textbf{Model}: LeNet-5 CNN~\cite{lecun1998gradient} trained from scratch (2 conv layers, 2 FC layers, $\sim$62K parameters).

\textbf{Training}: 10 rounds; 1 local epoch; batch size 64; SGD lr=0.01.

\subsubsection{Results}

\cref{fig:cnn_comparison} shows the accuracy of IID and non-IID with alpha value $\{0.30, 0.50\}$. For the IID setting, the accuracy of FiZK remains steady across all attacks and Byzantine ratios at 81.5--81.6\%. Multi-Krum struggles with IID model poisoning at 32.5\% when $\alpha = 0.30$ and fails when $\alpha = 0.50$ under label flip and model poisoning because the gradient norm explodes. Vanilla also collapses to 10\% when $\alpha = 0.30$ under label flip and model poisoning. For the non-IID setting when $\alpha = 0.30$, the accuracy of FiZK remains robust at 76.3--76.4\%, whereas Vanilla collapses to 10\% when $\alpha = 0.30$ under label flip and model poisoning. Multi-Krum struggles when $\alpha = 0.30$ under the backdoor attack at 69.8\%. For the non-IID setting when $\alpha = 0.30$, the accuracy of FiZK collapses when $\alpha = 0.30$ under label flip at 60.8\% and under model poisoning at 56.3\%. This demonstrates that the last problem is handling the majority Byzantine non-IID setting. This is because the ZK circuit currently relies on the sum of the gradients instead of the Euclidean norm. With the non-IID setting, the malicious clients are able to send gradients whose individual components sum to zero. The ZK-Firewall uses the sum of the gradients instead of the Euclidean norm. Hence, the malicious clients are able to send gradients whose individual components sum to zero. A verifiable training circuit that actually works i.e., that cryptographically proves that the gradient was produced by the actual forward and backward pass over the committed data would be the clincher, as a Byzantine attacker would be unable to pass such a circuit for any data distribution and any fraction of Byzantine clients. In that case, we would be able to claim that FiZK is unconditionally robust. Until that day, the existing sum-check proxy is a reasonable, albeit imperfect, approximation. We achieve perfect detection for all attacks with $\mathrm{TPR} = 100\%$, $\mathrm{FPR} = 0\%$, $\mathrm{F1} = 1.0$ for $\alpha = 0.30$ with IID. We also achieve perfect detection for all adaptive attacks, i.e., gradient substitution and compromised aggregator attacks.

\textbf{The Key Finding:} The ZK-Firewall is shown to generalize to CNNs with at least 81\% IID accuracy and perfect detection for all attacks, except for the majority Byzantine non-IID case with $\alpha = 0.50$.

\begin{figure*}[t]
  \centering
  \includegraphics[width=\textwidth]{cnn_full_comparison.pdf}
  \caption{CNN (LeNet-5/Fashion-MNIST) attack robustness: test accuracy (\%) for Vanilla FedAvg, Multi-Krum, and FiZK under Backdoor, Label Flip, and Model Poisoning. Top row: IID; bottom row: Dirichlet non-IID ($\beta=0.5$). FiZK consistently holds 81--82\% IID and 76\% non-IID at $\alpha=0.30$; Multi-Krum degrades to 10--32\% under poisoning attacks. \textbf{N/A}: Multi-Krum training failed at this configuration (gradient norms $\to\infty$; runs did not converge).}
  \label{fig:cnn_comparison}
\end{figure*}

%                     
\subsection{Cross-experiment analysis}
\label{sec:analysis}
%                     

\subsubsection{Threat coverage}

% We map each attack type to its corresponding defense layers in Table~1. 
% The ZK-Firewall (Layer~1, Stage~2) enforces gradient norm bounds before aggregation, blocking high-magnitude malicious updates. 
% The remaining layers (2--4) then mitigate smaller but potentially adversarial gradients that may still pass the initial check.

We can match each type of attack to the corresponding defense layer as indicated in \cref{tab:threats}. For the ZK-Firewall, the gradient norm is constrained before the aggregation (Layer 1, Stage 2). The remaining layers (Layers 2-4) will handle the small gradients that could be considered an attack but have already managed to evade the initial defense.

\begin{table}[t]
  \centering
  \small 
  \caption{Defense Layer Threat Coverage. We denote $\bullet\bullet$ as primary, $\circ$ as supporting, and -- as not applicable.}
  \label{tab:threats}
  \begin{tabular}{@{}lcccc@{}}
    \toprule
    \textbf{Attack} & \textbf{L1} & \textbf{L2} & \textbf{L3} & \textbf{L4} \\
    \midrule
    Gradient tampering   & $\bullet\bullet$ & --               & --               & --               \\
    Replay attack        & $\bullet\bullet$ & --               & --               & --               \\
    Exploding gradient   & $\bullet\bullet$ & --               & --               & --               \\
    Model poisoning      & $\circ$          & $\bullet\bullet$ & $\bullet\bullet$ & $\circ$          \\
    Label flipping       & --               & $\circ$          & $\circ$          & $\bullet\bullet$ \\
    Backdoor injection   & --               & $\circ$          & $\circ$          & $\bullet\bullet$ \\
    Gaussian noise       & $\bullet\bullet$ & $\circ$          & $\circ$          & --               \\
    Adaptive attack      & $\circ$          & $\circ$          & $\bullet\bullet$ & $\circ$          \\
    Aggregator tampering & $\bullet\bullet$ & --               & --               & --               \\
    \bottomrule
  \end{tabular}
\end{table}

\subsubsection{Complexity analysis}

The per-round costs for each part are given in table \cref{tab:complexity}. Multi-Krum is the big hitter, with a cost of $O(n^2 d/G)$ over $G$ parallel galaxies, while the IVC accumulator verification is much cheaper at $O(\log m)$.


\begin{table}[t]
  \centering
  \small
  \caption{Per-Round Complexity by Component ($n$: total clients, $G$: galaxies, $d$: model dimension, $m$: ZK-validated clients selected by Multi-Krum)}
  \label{tab:complexity}
  \begin{tabular}{@{}ll@{}}
    \toprule
    \textbf{Component} & \textbf{Complexity} \\
    \midrule
    Merkle tree construction (per galaxy) & $O(n/G)$ \\
    ZK proof generation (per client)      & $O(d)$ \\
    ZK-Firewall verification (total)      & $O(n)$ \\
    Statistical anomaly detection         & $O(n \cdot d)$ \\
    Multi-Krum (per galaxy)               & $O((n/G)^2 \cdot d)$ \\
    IVC folding (per galaxy)              & $O(m)$ \\
    Global Merkle verification            & $O(G\cdot log(G))$ \\
    Reputation update                     & $O(n)$ \\
    \bottomrule
  \end{tabular}
\end{table}



\begin{table}[t]
  \centering
  \small
  \caption{Computational Overhead (MNIST, IID, clean, 10 clients, 3 trials).}
  \label{tab:overhead}
  \begin{tabular}{@{}lccc@{}}
    \toprule
    \textbf{Defense} & \textbf{Avg round time (s)} & \textbf{Overhead} & \textbf{ZKP share} \\
    \midrule
    Vanilla FedAvg & $5.17\,\pm\,0.31$ & $1.00\times$ & -- \\
    FiZK           & $45.51\,\pm\,1.85$ & $8.80\times$ & 75.8\% \\
    \bottomrule
  \end{tabular}
\end{table}

\begin{table}[t]
  \centering
  \small
  \caption{Test Accuracy (\%) Under Adaptive Attack (MNIST, IID, 10 Clients, 10 Rounds)}
  \label{tab:accuracy}
  \begin{tabular}{@{}lccc@{}}
    \toprule
    \boldmath$\alpha$ & \textbf{Vanilla} & \textbf{Multi-Krum} & \textbf{FiZK} \\
    \midrule
    0.20 & 92.15 & 92.14          & 92.14 \\
    0.30 & 92.07 & 92.05          & 92.06 \\
    0.40 & 92.05 & 92.01          & 92.01 \\
    0.50 & 92.07 & \textbf{90.85} & \textbf{92.06} \\
    \bottomrule
  \end{tabular}
\end{table}

For any Byzantine-resilient federated learning protocol among $f$ faulty clients, the communication complexity must be at least linear ($\Omega(n)$). The per-client work of FiZK is $O(d + \log n)$, which covers both the gradient and the client-side Merkle proof. It matches the lower bound up to a multiplicative logarithmic factor. With the IVC folding of each galaxy, the client-side proof work $m$ is reduced to a single constant-size accumulator: $\pi_g^{\text{fold}}$. On the global side, the work of the Merkle verification is $O(G)$, and it does not depend on the number of clients in each galaxy.

\subsubsection{Detection summary}

\cref{fig:mnist_comparison,fig:cnn_comparison} illustrate the defenses used in these experiments. FiZK attains an accuracy of at least 92\% on MNIST and $\geq$81\% on CNN under all attacks when detection is perfect ($\mathrm{TPR} = 100\%$, $\mathrm{FPR} = 0\%$). Multi-Krum fails for $\alpha \geq 0.50$ on both datasets, as expected from its theoretical requirement that $f < n/2 - 2$.

% CNN Tables Referenced in Experiment 2
\begin{table}[t]
  \centering
  \small
  \caption{FiZK defense against adaptive attacks (LeNet-5/Fashion-MNIST,
    IID, 10 clients, $\alpha=0.30$, 3 trials).}
  \label{tab:cnn_adaptive}
  \setlength{\tabcolsep}{5pt}
  \begin{tabular}{@{}lcccc@{}}
    \toprule
    \textbf{Adaptive Attack} & \textbf{Acc. (\%)} & \textbf{TPR (\%)} & \textbf{F1} & \textbf{Time (s)} \\
    \midrule
    Gradient Substitution & 81.5\,$\pm$\,0.2 & 100.0 & 1.00 & 75.2 \\
    Compromised Aggregator & 81.6\,$\pm$\,0.3 & 100.0 & 1.00 & 66.7 \\
    \bottomrule
  \end{tabular}
\end{table}

\begin{table}[t]
  \centering
  \small
  \caption{Scalability (LeNet-5/Fashion-MNIST, clean, IID, 3 trials).
    ZKP Total = prove + verify time per round.}
  \label{tab:cnn_scale}
  \setlength{\tabcolsep}{5pt}
  \begin{tabular}{@{}lrrrr@{}}
    \toprule
    \textbf{Defense} & \textbf{Clients/} & \textbf{Acc.} & \textbf{Round} & \textbf{ZKP} \\
                     & \textbf{Galaxies} & (\%)          & \textbf{(s)}   & \textbf{(\%)} \\
    \midrule
    Vanilla  & 10/2  & 81.8\,$\pm$\,0.4 & 5.0   & --   \\
    Multi-Krum & 10/2  & 81.6\,$\pm$\,0.3 & 5.0  & --   \\
    \textbf{FiZK} & 10/2  & 81.4\,$\pm$\,0.4 & 95.7\,$\pm$\,12.4 & 71.6 \\
    \textbf{FiZK} & 30/6  & 69.5\,$\pm$\,0.4 & 230.7\,$\pm$\,3.9 & 75.1 \\
    \textbf{FiZK} & 50/10 & 62.0\,$\pm$\,2.3 & 379.9\,$\pm$\,7.1 & 75.8 \\
    \bottomrule
  \end{tabular}
\end{table}
% ============================================================
\section{Limitations and future work}
\label{sec:future}

\subsection{Limitations}

\paragraph{Computational overhead.}
The cost of generating ZK proofs is around $8.80\times$ wall-clock overhead over vanilla FedAvg (\cref{tab:overhead}). The proportion of ZKP steps to a total round is around 75.8\% on MNIST and 71\% to 76\% on CNN for Multi-Krum. There is an additional cost to Multi-Krum given by $O\!\left((n/G)^2 d\right)$, which is another constraint on galaxy sizes.

\paragraph{Sum-check proxy.}
The ZK-Firewall shows that scalar sums $S{i}^{(l)}$ are proven instead of the full Euclidean norm $\| \nabla w{i}\|_{2}.$
In wide layers, positive and negative values within the gradient vector cancel out, enabling both honest and Byzantine
sums to remain low regardless of the magnitude of the gradients.
In practice, ZK is most reliable for rejecting Byzantine attacks in bias layers, while weight layers are more vulnerable
to evasion under non-IID majority Byzantine scenarios.

\paragraph{Bound contamination threshold.}
$B_l$ is anchored at $Q_{0.25}$, staying in the honest regime, while $\alpha < 0.75$ (\cref{eq:bound}). Our results verify robustness at $\alpha = 0.60$; beyond $\alpha \approx 0.75$, Byzantine clients overwhelm $B_l$ at the 25th percentile. This is not due to ZKP, but the norm-bound proxy itself. A circuit that directly verifies gradient correctness eliminates this threshold.

\paragraph{Galaxy assignment and scale.}
Static round-robin assignment is suboptimal, and Merkle trees alone do not guarantee gradient privacy (secure aggregation is required). Evaluation beyond 50 clients and across larger datasets remains future work.

\subsection{Future work}

\paragraph{Full verifiable training circuit.}
The main open direction here is the replacement of the layer sum proxy with a circuit that directly verifies $\nabla w_i = \nabla \mathcal{L}(w, D_i)$. Each SGD step now corresponds to one ProtoGalaxy folding step, and the final accumulator verifies the entire local training loop in $O(1)$ time. This would remove the majority Byzantine threshold entirely because no poisoned gradient will verify under the forward-backward pass constraint. SNARKs for neural network training are an active area of research in zkML.

\paragraph{Efficiency, privacy, and scale.}
Using GPU acceleration in the proving process (\eg, Icicle, cuZK) would reduce the time per client in the proving process ($\approx$1.8s per client on MNIST, $\approx$7s per client on CNN). Integrating FiZK with differential privacy or homomorphic encryption would introduce client gradient confidentiality. Global folding of galaxy proofs $\{\pi_g^{\mathsf{fold}}\}$ in a single accumulator using a second pass of ProtoGalaxy would reduce the time complexity of global server verification to $O(1)$. Finally, experimenting with larger datasets (\eg, CIFAR-10, medical imaging) with 100.

% ============================================================
\section{Conclusion}
\label{sec:conc}
% ============================================================

We proposed FiZK, a hierarchical federated learning framework combining zero-knowledge gradient validation, Merkle tree commitments, Byzantine-robust aggregation, and IVC-based proof folding. The ZK-Firewall enforces per-layer gradient norm bounds and sum correctness before any update enters aggregation, making Byzantine rejections cryptographically non-repudiable. Galaxy-level ProtoGalaxy folding compresses all per-client proofs into a constant-size accumulator $\pi_g^{\mathsf{fold}}$, reducing server-side verification to $O(1)$ per galaxy.

As shown in \cref{fig:mnist_comparison,fig:cnn_comparison}, FiZK consistently outperforms both baselines across IID and non-IID settings. On MNIST, FiZK maintains $\geq$92\% IID accuracy under model poisoning at $\alpha=0.60$ (majority-Byzantine), while Multi-Krum collapses to $\leq$11\% at $\alpha\geq0.50$; under non-IID conditions FiZK holds 87--90\% at $\alpha=0.30$. On Fashion-MNIST CNN, FiZK achieves perfect detection (TPR\,=\,100\%, FPR\,=\,0\%, F1\,=\,1.0) across all attack types under IID at $\alpha=0.30$, maintaining 81--82\% accuracy, while Multi-Krum degrades severely (32.5\,$\pm$\,29.3\% under IID model poisoning). The $8.80\times$ computational overhead is dominated by CPU-bound ZKP operations (75.8\% of round time) with communication cost identical to vanilla FedAvg. Overall, FiZK provides a practical and provably sound solution for verifiable Byzantine-robust federated learning, generalising across model architectures and data distributions.

{
    \small
    \bibliographystyle{ieeenat_fullname}
    \bibliography{main}
}

% WARNING: do not forget to delete the supplementary pages from your submission 
% \input{sec/X_suppl}

% Appendix removed for CVPR workshop submission (page limit).
% Algorithm and notation table are available in the full technical report.

\end{document}
